\section{An explicit frame-slab kernel}
\label{sec:sr-frame-slab-kernel}

We prove the archimedean integrability and reproduction used in
\eqref{eq:sr-slab-reproduction-input} for an explicit kernel.
This removes a separate analytic reproduction assumption once the
doubled Gaussian has been identified with this kernel in the chosen
highest-weight coordinates.

Put \(r=p+q\), \(J=\operatorname{diag}(1_p,-1_q)\), and
\[
 \mathcal D=\{Y\in\operatorname{GL}_r(\mathbb C):
             (Y^*JY)_{[1,p],[1,p]}>0\},\qquad U=U(p,q).
\]
Write \(\Delta_k(Y)\) for the leading \(k\)-minor,
\(\Delta_0=1\).  Let \(\lambda=(\lambda_1,\ldots,\lambda_r)\)
be integral, nonincreasing separately on the two blocks, with
\begin{equation}\label{eq:sr-classical-weight-gap}
 \lambda_i+n\le\lambda_j
           \quad(1\le i\le p<j\le r).
\end{equation}
These are the classical frame weights: the first block lies below
the second block.  They are not the algebraic weight array before
the dictionary to classical frames.

Define
\[
 m_\lambda(Y)=
    \det(Y)^{\lambda_r}
       \prod_{k=1}^{r-1}
          \Delta_k(Y)^{\lambda_k-\lambda_{k+1}},\qquad
 K_\lambda(Y,Z)=m_\lambda(JZ^*JY).
\]
Only the \(p\)-th minor can have a negative exponent in the
product; that minor is nonzero on \(\mathcal D\).
Indeed the first \(p\) columns of \(Y\) form a maximal positive
subspace, and projection to the first \(p\) coordinates is
injective on any such subspace.  Determinants are also nonzero.
Thus \(m_\lambda\) is holomorphic on \(\mathcal D\).

The same assertion for the cross matrix in \(K_\lambda\) follows
by writing the two positive \(p\)-planes as graphs
\((I,A)^t,(I,B)^t\) with \(\|A\|,\|B\|<1\).
Their cross pairing has the invertible factor
\(I-B^*A\), since \(\|B^*A\|<1\).
Consequently \(K_\lambda\) is holomorphic in \(Y\) and
antiholomorphic in \(Z\).
For an invertible upper triangular matrix \(b\), put
\(\chi_\lambda(b)=\prod_i b_{ii}^{\lambda_i}\).
Leading minors give, directly,
\begin{align}
 K_\lambda(Yb,Z)&=\chi_\lambda(b)K_\lambda(Y,Z),\label{eq:sr-kernel-weight}\\
 K_\lambda(Y,Zb)&=\overline{\chi_\lambda(b)}K_\lambda(Y,Z),\nonumber\\
 K_\lambda(uY,uZ)&=K_\lambda(Y,Z)\quad(u\in U),\nonumber\\
 K_\lambda(Z,Y)&=\overline{K_\lambda(Y,Z)},\qquad
 K_\lambda(I,I)=1.\nonumber
\end{align}
For the second weight identity, \(Jb^*J\) is lower triangular,
so its leading \(k\)-minor multiplies that of the cross matrix
by \(\prod_{i\le k}\overline{b_{ii}}\).
For symmetry use
\(JY^*JZ=J(JZ^*JY)^*J\); the two diagonal \(J\)'s cancel in
each leading minor.

\begin{srlemma}[Lowest weight uniqueness]
\label{lem:sr-frame-lowest-unique}
Suppose \(F\) is holomorphic on \(\mathcal D\), has the right
upper-triangular character \(\chi_\lambda\), and for every
diagonal unitary \(t\) satisfies
\(F(tY)=\chi_\lambda(t)F(Y)\).
Then \(F(Y)=F(I)m_\lambda(Y)\) on \(\mathcal D\).
\end{srlemma}

\begin{proof}
Both \(m_\lambda\) and \(F\) have these two transformation laws.
Their difference \(G=F-F(I)m_\lambda\) vanishes at \(I\).

First take an invertible lower triangular \(Y\in\mathcal D\) and
let \(D(z)=\operatorname{diag}(1,z,\ldots,z^{r-1})\).
The matrix
\(Y(z)=D(z)YD(z)^{-1}\), initially for \(z\ne0\), extends
polynomially to \(z=0\), since its entries below the diagonal
are \(z^{i-j}Y_{ij}\).  Its value at zero is the diagonal of \(Y\).
For \(0<|z|\le1\) it stays in \(\mathcal D\):
before the invertible right diagonal change, the positive
coordinates are scaled by at least \(|z|^{p-1}\), and the
negative coordinates by at most \(|z|^p\).
This preserves strict positivity of the first \(p\)-plane.
At \(z=0\) the diagonal matrix is also in \(\mathcal D\).
On \(|z|=1\), the two transformation laws cancel to give
\(G(Y(z))=G(Y)\).
The holomorphic function of the disk therefore has this constant
boundary value and is constant, by the maximum principle.
At zero its value is
\(\chi_\lambda(\operatorname{diag}Y)G(I)=0\).
Hence \(G(Y)=0\) on these lower triangular matrices.

On the open cell where every leading minor is nonzero,
Gaussian elimination writes a matrix as a lower triangular
matrix times an upper triangular matrix, so the right
transformation law proves \(G=0\) there.
That cell is dense in \(\mathcal D\):
replace \(Y\) by \(Y+sI\), avoiding the finite root set of each
of its monic leading-minor polynomials in \(s\).
Continuity proves \(G=0\) everywhere.
\end{proof}

\begin{srlemma}[The unitary minor bound]
\label{lem:sr-unitary-minors}
For \(u\in U(p,q)\),
\[
 1\le|\Delta_p(u)|,\qquad |\det u|=1,\qquad
 |\Delta_k(u)|\le|\Delta_p(u)|\quad(0\le k\le r).
\]
Consequently
\begin{equation}\label{eq:sr-frame-kernel-decay}
 |m_\lambda(u)|\le|\Delta_p(u)|^{-n}.
\end{equation}
If \(p=0\) or \(q=0\), interpret the last bound as \(1\).
\end{srlemma}

\begin{proof}
In block form \(u=\left(\begin{smallmatrix}A&B\\ C&D\end{smallmatrix}\right)\),
the relation \(u^*Ju=J\) gives \(A^*A-C^*C=I\).
Thus all singular values of \(A\) are at least one.
The determinant identity gives \(|\det u|=1\).

For the remaining minor inequality let \(R_p\) have the first
\(p\) rows of \(u\) and the last \(q\) rows of \(I\);
let \(R'_p\) have the first \(p\) rows of \(I\) and the last
\(q\) rows of \(u\).
The unitary relation gives
\((R'_p)^*R'_p=R_p^*R_p\), so \(V=R'_pR_p^{-1}\) is unitary.
Let \(R_k\) similarly use the first \(k\) rows of \(u\) and
the remaining rows of \(I\).
Every row of \(R_kR_p^{-1}\) is a row of either \(I\) or \(V\),
and therefore has norm one.
Hadamard's determinant inequality gives
\[
 |\Delta_k(u)/\Delta_p(u)|
       =|\det(R_kR_p^{-1})|\le1.
\]
In \(m_\lambda(u)\), every exponent
\(\lambda_k-\lambda_{k+1}\) except that at \(k=p\)
is nonnegative; bound its minor by \(|\Delta_p(u)|\).
The \(p\)-th factor is already a power of that minor, while
the full determinant has modulus one.
The resulting exponent is
\(\lambda_1-\lambda_r\le-n\), by
\eqref{eq:sr-classical-weight-gap}.
This proves \eqref{eq:sr-frame-kernel-decay}.
The definite cases use the same argument with all leading minors
bounded by one and the full determinant of modulus one.
\end{proof}

Let
\[
 S_\delta=\{Y:\ |(Y^*JY-J)_{ab}|<\delta\quad(1\le a,b\le r)\}.
\]
For sufficiently small \(\delta>0\), this open set lies in
\(\mathcal D\) and contains \(I\).
Indefinite Gram--Schmidt writes every member uniquely as
\(Y=ub\), where \(u\in U(p,q)\), \(b\) is upper triangular
with positive real diagonal, and \(b^*Jb=Y^*JY\).
All these \(b\)'s lie in a compact neighborhood of \(I\).
To see the uniform bound, the successive pivots in the
indefinite Cholesky recursion are near \(+1\) for the first block
and \(-1\) for the second; they remain bounded away from zero.
Every division and positive square root in the recursion is
therefore bounded on the closure of a sufficiently small
neighborhood of \(J\).

Ambient Lebesgue measure in these coordinates is
\(j(b)\,du\,db\), with \(j\) smooth and bounded on this compact
set of \(b\)'s.  This follows either from the change-of-variables
Jacobian or from its local principal-bundle coordinates.
There is no \(u\)-dependent factor: left multiplication by
\(u\) preserves ambient Lebesgue measure, since
\(|\det_\mathbb C u|^{2r}=1\), and preserves Haar measure \(du\).

\begin{srproposition}[Slab integrability]
\label{prop:sr-frame-slab-integrable}
If \(n\ge2r\), then for sufficiently small \(\delta>0\),
\(K_\lambda(X,\cdot)\) is bounded and absolutely integrable
on \(S_\delta\), for every \(X\in\mathcal D\).
\end{srproposition}

\begin{proof}
First take \(X=I\).  The weight relation, the compact bound on
\(b\), and Lemma~\ref{lem:sr-unitary-minors} reduce the integral to
\(\int_U|\Delta_p(u)|^{-n}\,du\).
Use the Cartan integration formula with
\(\ell=\min(p,q)\), \(t_1\ge\cdots\ge t_\ell\ge0\).
The positive-block determinant has modulus
\(\prod_i\cosh t_i\), and, up to a fixed constant, the radial
Jacobian is
\[
 \prod_{i<j}\sinh^2(t_i-t_j)\sinh^2(t_i+t_j)
 \prod_i\sinh^{2|p-q|}(t_i)\sinh(2t_i).
\]
Here is the Jacobian calculation.  Ordinary singular-value
decomposition of the off-diagonal block, followed by
\(A^*A-C^*C=I\) and its companion identities, puts \(u\), up
to left and right multiplication by \(U(p)\times U(q)\), into
\(\ell\) hyperbolic planes
\(\left(\begin{smallmatrix}\cosh t_i&\sinh t_i\\
\sinh t_i&\cosh t_i\end{smallmatrix}\right)\), with identity
on the unused coordinates.
On the regular part, differentiate this compact--diagonal--compact
map and divide out the compact centralizer.
For each pair of opposite restricted root spaces the two
compact directions are sent to their difference under
\(\exp(\operatorname{ad}H)\); its nonconstant determinant factor
is \(\sinh(\alpha(H))\) in each real root direction.
Matrix units coupling two hyperbolic planes give
\(e_i-e_j,e_i+e_j\), each with real multiplicity \(2\);
coupling a plane to an unused coordinate gives \(e_i\)
with multiplicity \(2|p-q|\); and the remaining imaginary
direction within a plane gives \(2e_i\) with multiplicity \(1\).
Their product is exactly the displayed Jacobian, up to a
positive constant.  The compact factors have finite volume.
The singular walls have lower dimension and contribute measure
zero; restricting to the ordered chamber only changes the
constant by the finite Weyl multiplicity.
Bounding \(\sinh t\le e^t\) and
\(\cosh t\ge e^t/2\), its integrand is at most a constant times
\[
 \exp\left[-\sum_{i=1}^{\ell}
                   (n-2r+4i-2)t_i\right]
 \le \exp\left[-2\sum_i t_i\right].
\]
The integral converges.  The same bounds and
\(|\Delta_p(u)|\ge1\) prove boundedness.
When \(\ell=0\), \(U\) is compact and the assertion is immediate.

For arbitrary \(X\), indefinite Gram--Schmidt gives
\(X=u_0b_0\).  Kernel invariance and its first weight relation
reduce \(K_\lambda(X,Y)\) to a constant times
\(K_\lambda(I,u_0^{-1}Y)\).
The slab and Lebesgue measure are invariant under \(u_0\).
This proves both conclusions.
\end{proof}

\begin{srproposition}[Slab reproduction]
\label{prop:sr-frame-slab-reproducing}
Under the same conditions there is \(\rho_\lambda>0\) such that
for every holomorphic \(F\) on \(\mathcal D\) with right
character \(\chi_\lambda\) and
\(\sup_{u\in U}|F(u)|<\infty\),
\begin{equation}\label{eq:sr-explicit-slab-reproduction}
 \int_{S_\delta}K_\lambda(X,Y)F(Y)\,dY
                 =\rho_\lambda F(X)\qquad(X\in\mathcal D).
\end{equation}
\end{srproposition}

\begin{proof}
The factorization \(Y=ub\) and compactness of the possible \(b\)'s
make \(F\) bounded on the slab.
All following integrals are therefore absolutely convergent.

Let \(T\subset U\) be the diagonal compact torus with probability
Haar measure and set
\[
 G(Y)=\int_T\chi_\lambda(t)^{-1}F(tY)\,dt .
\]
Compact averaging preserves holomorphy and the right character.
It gives the left torus character in
Lemma~\ref{lem:sr-frame-lowest-unique}, and \(G(I)=F(I)\).
Thus \(G(Y)=F(I)m_\lambda(Y)=F(I)K_\lambda(Y,I)\).

Kernel invariance and its weight at the base point show
\(K_\lambda(I,tY)=\chi_\lambda(t)^{-1}K_\lambda(I,Y)\).
Changing variables in the slab, which \(T\) preserves, consequently
gives
\[
 \int_{S_\delta}K_\lambda(I,Y)G(Y)\,dY
       =\int_{S_\delta}K_\lambda(I,Y)F(Y)\,dY .
\]
Using the formula for \(G\) proves reproduction at \(I\), with
\[
 \rho_\lambda=\int_{S_\delta}|K_\lambda(I,Y)|^2\,dY .
\]
This is finite by boundedness and integrability of the kernel.
It is positive because the slab is open, contains \(I\), and the
continuous kernel equals one at \((I,I)\).
For \(X=ub\), change variables by \(u\), then use the first
weight identity to reduce to the identity at \(I\) for
\(Y\mapsto F(uY)\).  This proves
\eqref{eq:sr-explicit-slab-reproduction}.
\end{proof}

Products over the archimedean embeddings have the same proof:
take the product of the domains, tori, slabs and kernels.
Their reproduction constant is the positive product of the
individual \(\rho_\lambda\)'s.

This also gives a short and exact test for the doubled Gaussian
normalization.  If the paired-derivative kernel \(K_{\rm Gau}\)
is holomorphic/antiholomorphic, unitary invariant, and has the
two characters in \eqref{eq:sr-kernel-weight}, then
Lemma~\ref{lem:sr-frame-lowest-unique} gives
\[
 K_{\rm Gau}(Y,I)=K_{\rm Gau}(I,I)m_\lambda(Y).
\]
Writing the second frame as \(ub\) proves
\(K_{\rm Gau}=K_{\rm Gau}(I,I)K_\lambda\) everywhere.
In orthonormal Gaussian coordinates its diagonal is computed
by
\[
 P(\partial_z)\overline P(\partial_w)
          \exp\!\left(\sum_j c_jz_jw_j\right)\big|_{z=w=0}
      =\sum_\alpha\alpha!c^\alpha|P_\alpha|^2,
      \qquad c_j>0 .
\]
For a nonzero orthonormal harmonic family this is strictly
positive.  Hence the normalized paired Gaussian has a positive
scalar multiple of the explicit kernel, and reproduction has
that same positive scalar multiplier.
For the actual theta model, the three stated structural
identities and the sign of the Gaussian coordinates must be
verified in its fixed frame and differential dictionary.
The argument above proves all the ensuing slab convergence,
positivity and reproduction, with no further global
automorphic input.

\section{A common theta denominator at all finite places}
\label{sec:sr-theta-common-denominator}

The denominator in the short theta spanning construction can be chosen
before its auxiliary lattice, Fourier label, harmonic polynomial, or
section.  We give its exact formula and the three valuation arguments.
The argument explicitly records the local geometric calculus it uses;
establishing that calculus for the chosen integral models remains part
of their construction.

Let \(E\) be the fixed Galois CM host, put
\[
 f=[E:\mathbb Q]/2,\quad r=p+q,\quad
 \Delta=|\operatorname{disc}E|,\quad
 d_P=\sum_{\tau,i}a_{\tau,i},\quad
 \kappa=d_P+frn/2,
\]
and let \(M\ge3\) be the fixed odd level.
Here \(n\) is even, and \(P\) ranges over the prescribed integral
harmonic polynomials of row degrees \(a_{\tau,i}\); every nonzero
such \(P\) is homogeneous of total degree \(d_P\).
The finite Fourier-label group has order
\[
 Q_M=M^{frn}.
\]
Take the actual Fourier datum of a level class \(x_0\):
its coefficients are
\(w_{x_0}(q)=\zeta_M^{\langle x_0,q\rangle}\).
They satisfy \(w_{x_0}(q)^M=1\), so they are algebraic integers
at every finite place, including places dividing \(M\).

The single integer used below is
\begin{equation}\label{eq:sr-common-theta-denominator}
 Z=Q_M\,\Delta^{d_P}M^{d_P}
   =M^{frn+d_P}\Delta^{d_P},\qquad N_a=Z^2.
\end{equation}
Its independence from the theta vector is essential.

Let \(\Lambda_a\) be an auxiliary positive definite unimodular
Hermitian lattice and
\(\Lambda_n=\Lambda_a\cap\mathcal O_E^n\).
For each embedding \(\tau:E\hookrightarrow K\), let
\(\det_\tau\Lambda_a,\det_\tau\Lambda_n\) be the determinant
fractional ideals in the common coefficient field \(K\).
For the fixed signature multiplicities \(r_\tau\), write
\(\mathcal T(J)=\prod_\tau\tau(J)^{r_\tau}\).
The scalar \(s\in K\) is chosen with the exact balancing condition
\begin{equation}\label{eq:sr-reflex-scalar-balance}
 (s)\mathcal T(\det\Lambda_n)
          \subset\mathcal T(\det\Lambda_a).
\end{equation}
This is an ideal condition; it is not replaced by an archimedean
norm estimate.

We use the following local reading of the algebraic theta section
\(Y_{\Lambda,P,x_0}\).
At a divisorial integral chart, after a finite extension of complete
DVRs, it is
\begin{equation}\label{eq:sr-theta-local-reading}
 Y_{\Lambda,P,x_0}
    =Q_M^{-1}\sum_{q}w_{x_0}(q)H_q,\qquad
 H_q=I_\Lambda^{-1}N^{-d_P}
       \sum_{z\in\mathcal Z_\Lambda}
                {\cal D}_P(F_{z,q}).
\end{equation}
Here \(N\) is an odd auxiliary denominator with \(M\mid N\),
\(I_\Lambda\) is the prescribed lattice index,
\(\mathcal Z_\Lambda\) is its isotropic torsion group, and the
\(F_{z,q}\) are the normalized symmetric translates of the
normalized multiplication pullback of the same primitive theta
section.  The differential \({\cal D}_P\) is that of the fixed
canonical connection and frame.
The theta scalar multiplying the source's short rule is included
in \(Y\); it is not an extra unit inserted after taking valuations.
Formula \eqref{eq:sr-theta-local-reading} is an equality in the
fraction field, valid with the same field embedding as \(s,P\).

Normalize a chart valuation as \(v=v_T/e\).
Under a finite extension of ramification index \(t\), use
\(v_{T'}/(te)\).  Its restriction is exactly \(v_T/e\).
Thus an inequality established after the extension descends
without a new loss.
The required connection bound, for every integral local germ of
the indicated weight and every line power, has loss
\begin{equation}\label{eq:sr-chart-ramification-loss}
 \rho(v)=\frac{d_P}{2}
   \sum_{\ell\mid\operatorname{cond}(E)}
       v(\ell)\frac{\operatorname{ord}_\ell\Delta}{[E:\mathbb Q]}
 \le d_P\,v(\Delta).
\end{equation}
All terms are nonnegative, and
\(v(\Delta)=\sum_\ell\operatorname{ord}_\ell(\Delta)v(\ell)\);
this proves the last inequality.
The more precise half-degree loss on the left is the one in the
ramified CM connection estimate.

We first explain two elementary pieces of the local argument.
For a homogeneous harmonic \(P\), normalized pullback along a CM
endomorphism \(\beta\) obeys
\begin{equation}\label{eq:sr-normalized-harmonic-pullback}
 {\cal D}_P([\beta]^*_{\rm norm}h)
        =\operatorname{Nm}_a(\beta)\,{\cal D}_P(h),
 \quad
 \operatorname{Nm}_a(\beta)
   =\prod_{\tau,i}\tau(\beta)^{a_{\tau,i}},
 \quad
 \operatorname{Nm}_a(m)=m^{d_P}.
\end{equation}
The formula uses the actual normalized pullback, its canonical
connection, and its induced integral tangent directions.
It follows by applying the differential chain rule in each
eigendirection: every letter in row \((\tau,i)\) gains
\(\tau(\beta)\), and every monomial has the same row degrees.
This also shows multiplicativity of \(\operatorname{Nm}_a\).

Suppose \(N=\beta\alpha n'\) in \(\mathcal O_E\), with \(n'\)
a unit at the chart.  If the translated section at level \(N^2\)
is the normalized \(\beta\)-pullback of a regular reached section
\(h\), then \eqref{eq:sr-normalized-harmonic-pullback} gives
\[
 \frac{\operatorname{Nm}_a(\alpha)}{N^{d_P}}
             {\cal D}_P(F)
     =(n')^{-d_P}{\cal D}_P(h).
\]
Regularity puts the germ of \(h\) in the integral power-series
ring.  The canonical connection estimate
\eqref{eq:sr-chart-ramification-loss} therefore gives valuation
at least \(-\rho(v)\); the factor \(n'\) has valuation zero.
The equality expressing the mixed symmetric translate as that
normalized pullback, and regularity of the reached section,
are the geometric assertions needed at this point.

There is also an exact cancellation of primary torsion.
Write \(\mathcal Z_\Lambda=A_\ell\times A^{(\ell)}\).
When \(\Lambda_a\) is standard at \(\ell\), every element of
\(A_\ell\) is killed by \(N\), and
\[
 |A_\ell|=\ell^{\operatorname{ord}_\ell I_\Lambda},
 \qquad I_\Lambda=|A_\ell|I^{(\ell)},\qquad
 \ell\nmid I^{(\ell)}.
\]
The normalized multiplication pullback \([N]^*_{\rm norm}\theta\)
is fixed by its normalized symmetric \(N\)-torsion translations.
Translations of coprime orders commute: their commutator is a
Weil-pairing value of order dividing both orders, hence equals
one.  Their product is the normalized translate of their product
point, by uniqueness of the symmetric normalization at odd level.
It follows that \(F_{a,a',q}\) is independent of \(a\in A_\ell\)
once \(a'\) is fixed.
This is exact equality of sections, including their phases.
Thus
\begin{equation}\label{eq:sr-primary-index-cancellation}
 I_\Lambda^{-1}\sum_{a,a'}{\cal D}_P(F_{a,a',q})
       =(I^{(\ell)})^{-1}\sum_{a'}{\cal D}_P(F_{0,a',q}).
\end{equation}
The notation \srcid{sqN} in the frozen source only identifies
the equal exponents \((N:\mathbb Z)^2\) and \(N^2:\mathbb N\).
No squaring of a section is used to obtain this cancellation.

Here are the three cases for the local reading.
They use the same connection, the same normalized translations,
and the same scalar condition \eqref{eq:sr-reflex-scalar-balance}.

\paragraph{Case \(\ell\nmid M\).}
The level translates have prime-to-\(\ell\) order.
The adapted lattice isogeny separates the \(\ell\)-primary
isotropic sum and its prime-to-\(\ell\) part.
The normalized isogeny trace cancels the
\(\ell\)-primary part of \(I_\Lambda\).
For the remaining integral harmonic germ its determinant factor
is
\(\prod_\tau\tau(\det A_\ell)^{r_\tau}\), where \(A_\ell\)
is the integral adapted lattice matrix.
The balancing condition makes
\begin{equation}\label{eq:sr-adapted-determinant-integral}
 v\left(s\prod_\tau
                  \tau(\det A_\ell)^{r_\tau}\right)\ge0.
\end{equation}
To verify the ideal step explicitly, let \(G\) be a local column
basis of \(\Lambda_n\), and let \(g\) carry
\(\Lambda_a\) to the standard lattice at \(\ell\).
The adapted matrix is a rational-integer unit times \(gG\);
the multiplying integer is prime to \(\ell\).
Hence
\[
 v(\mathcal T(\det\Lambda_n))
       =\sum_\tau r_\tau v(\tau(\det G)),\quad
 v(\mathcal T(\det\Lambda_a))
       =-\sum_\tau r_\tau v(\tau(\det g)).
\]
Applying \(v\) to \eqref{eq:sr-reflex-scalar-balance} proves
\eqref{eq:sr-adapted-determinant-integral}.
Using it in the normalized isogeny trace and applying
\eqref{eq:sr-normalized-harmonic-pullback} gives
\[
                         v(sH_q)\ge-\rho(v).
\]
The existence of the adapted isogeny and the normalized trace
identity, with this determinant factor, are required here.

\paragraph{Case \(\ell\mid M\), \(\ell^2\nmid M\).}
Take the rational integer \(\alpha=\ell\), so
\(\operatorname{Nm}_a(\alpha)=\ell^{d_P}\).
The odd-level torsion factorization expresses the mixed
translate as a normalized pullback of a reached section with
the residual auxiliary denominator prime to \(\ell\).
The preceding chain-rule argument and primary decomposition
therefore give
\[
                         v(\ell^{d_P}H_q)\ge-\rho(v).
\]
Write \(M=\ell M'\); \(M'\) is integral, so multiplication by
\((M')^{d_P}\) preserves the inequality.
Since the auxiliary lattice is standard at every prime of \(M\),
its two determinant ideals in
\eqref{eq:sr-reflex-scalar-balance} are units at the chart.
Thus \(v(s)\ge0\), and
\[
                         v(M^{d_P}sH_q)\ge-\rho(v).
\]
This uses the rational \(\ell\) as the CM endomorphism.
There is no need to choose an element of norm \(\ell\).

\paragraph{Case \(\ell^2\mid M\).}
Write
\[
 N=\ell^e n',\qquad M=\ell^k m,\qquad
 k\le e,\quad m\mid n',\quad \ell\nmid n',
\]
and take \(\alpha=\ell^k\), \(\beta=\ell^{e-k}\).
Then
\(\alpha\overline\alpha=\ell^{2k}\),
\(\beta\overline\beta=\ell^{2(e-k)}\).
A translated point \(t\) satisfies \(t^{NM}=1\), because
\(t^N\) is killed by the level ideal and hence by \(M\).
Consequently
\[
 t^{N^2}=1,\qquad
 (t^\beta)^{(n')^2\ell^{2k}}=1:
\]
the first exponent is a multiple of \(NM\), and the second is
\(\ell^{e+k}(n')^2=(NM)(n'/m)\).
These are precisely the torsion exponents needed to express the
mixed translate as a normalized \(\beta\)-pullback of a regular
reached section at line power \((n')^2\ell^{2k}\).

Use \eqref{eq:sr-primary-index-cancellation}; its remaining index
\(I^{(\ell)}\) is a unit.
Apply the pullback calculation with
\(\operatorname{Nm}_a(\alpha)=\ell^{kd_P}\).
Every remaining summand has valuation at least \(-\rho(v)\),
so the nonarchimedean triangle inequality gives
\[
                   v(\ell^{kd_P}H_q)\ge-\rho(v).
\]
The same determinant-ideal argument gives \(v(s)\ge0\), and
\(M/\ell^k=m\) is integral.  Thus again
\[
                         v(M^{d_P}sH_q)\ge-\rho(v).
\]

In each of the three cases, the integral Fourier coefficients and
\eqref{eq:sr-theta-local-reading} imply
\[
 v(sY_{\Lambda,P,x_0})
      \ge-v(Q_M)-d_Pv(M)-\rho(v).
\]
Multiplication by \(Z^2\), followed by
\eqref{eq:sr-chart-ramification-loss}, gives a nonnegative valuation.
The normalization of valuations under finite DVR extension noted
above descends it to the original divisorial chart.
If the original section already belongs to the horizontal cusp
space, this proves membership in the prescribed integral cusp
lattice, whose additional conditions are exactly these
nonnegative vertical valuations.
The same \(N_a=Z^2\) works for every auxiliary lattice, every label,
and every polynomial of the fixed row degrees.

Finally the size estimate is elementary and has no local choice
in it.  Suppose
\[
 \log M\le30L,\qquad
 \log\Delta\le7fL,\qquad f,r,L\ge1,\qquad z=frL .
\]
Since \(frn+d_P\le2\kappa\) and \(d_P\le\kappa\),
\[
 \log Z
   =(frn+d_P)\log M+d_P\log\Delta
   \le60L\kappa+7fL\kappa
   \le67z^2\kappa ,
 \qquad
 \log N_a\le134z^2\kappa .
\]
Thus the integer clearing factor contributes one fixed exponent
before the theta spanning family is chosen.

The denominator size, determinant-ideal implications, primary
factorization, and valuation bookkeeping above are complete.
For their application to the actual arithmetic models the local
reading \eqref{eq:sr-theta-local-reading}, the half-degree
connection bound \eqref{eq:sr-chart-ramification-loss}, regularity
and exact normalized pullback of reached sections, and the adapted
isogeny trace with the displayed determinant factor must be
established in those models.
These are stronger than integrality of the Fourier coefficients.
Their producers in the frozen source include
\srcid{HarmonicValuationLowerBoundHostFieldAllPlaces},
\srcid{HarmonicValueTranslatedProductValuationBound},
\srcid{ReachedSectionRegularOnIntegralModel},
\srcid{MixedTranslatePullbackEqReachedSection},
\srcid{ExistsAdaptedSimilitudeDataAtPrime},
\srcid{SumHvPrimaryKernel}, and
\srcid{ShortSeriesBoundCoprimeTorsionMul}.
The global assembly is
\srcid{ExistsBoundedIntegerGoodCoverThetaSeriesIntegral}.
